\scanpage{14}
% 1–6
Soient $X$ un schéma $A_*$-formel et $Y$ un sous-schéma formel (lisse) de $X$. \deleted{On dit} On dit que $Y$ est un \deleted{\unsure{sous}} $A_*$-schéma $A_*$-formel s'il existe sur $Y$ une $A_*$ structure telle que l'inclusion can. soit un $A_*$-morphisme. Cette $A_*$-structure est alors unique.

% 6–14
En effet, il est immédiat qu'on peut (sans hyp. sur $Y$) trouver un \inserted{$A_*$} morphisme $X\to Y'$ ($Y'$ un schéma $A_*$-formel) tel que le morphisme induit $\deleted{X}Y\to Y'$ soit un isomorphisme formel. \deleted{\unsure{Considérons|Cons.}} Considérons l'isomorphisme inverse, et le composé
\diagram{\begin{tikzcd}[ampersand replacement=\&]
\& Y\arrow[dr] \& \\
Y'\arrow[ur,"\sim"]\arrow[rr] \& \& X
\end{tikzcd}}.
Utilisant (\unsure{e|c|d}), on voit que \deleted{$\unsure{A_*\text{ est}}$} $Y$ est un sous-schéma $A_*$-formel ssi $\deleted{\unsure{X\to}}Y'\to X$ est un $A_*$-morphisme ; et alors la $A_*$-structure de $Y$ est celle déduite.

% 15–17
\uline{Exemples} La diagonale de $X\times X$ est un sous $A_*$-schéma formel. Plus gén. celle de $X\times_Y X$ si $X\to Y$ est un $A_*$ morphisme submersif \deleted{(i.e. \unsure{formt. lisse|form. lisse|formel lisse})}

% 16–20, left
% slanted
\margin{Un morphisme formel $X\to Y$ entre \unsure{schs|objts} $A_*$-formels est un $A_*$-morphisme ssi son graphe dans $X\times Y$ est un sous-schéma $A_*$-formel}

% 18–27
Soit $\mathfrak X$ un schéma sur $\mathcal O_{A_*}(X)$ ($X$ un $A_*$-schéma formel) muni d'une section \inserted{$s_0$} sur $\mathcal O_{A_*}(X)$, supposons que $\mathfrak X$ soit lisse sur $\mathrm{Spec}\,\mathcal O_{A_*}(\deleted{X}X)$ aux points de cette section. Considérons le complété formel \inserted{$\widehat{\mathfrak X}\deleted{\unsure{^{s_0}}}$} de $\mathfrak X\otimes_{\mathcal O_{A_*}(X)}\mathcal O(X)$ le long de la section $s$ déduite de $s_0$. Alors une construction utilisant e) permet d'y mettre une \inserted{$A_*$-}structure, \inserted{telle que $X\to\widehat{\mathfrak X}\to X$ \unsure{soit|soient} \deleted{\unsure{de|des}} $A_*$-morphisme.} \deleted{\unsure{le|la}} $\mathfrak X$ \deleted{\unsure{est}} dépend $A_*$-fonctoriellement du triple $(X,\mathfrak X,s_0)$, sa formation est compatible avec les images inverses (par $X'\to X$\ldots)
% archive sheet number: 13
